\chapter[Sermon 84. The Fact of St. John Is Declared, Which Would Have Worshipped the Angel, and of the Angel Prohibiting.]{The Fact of St. John Is Declared, Which Would Have Worshipped the Angel, and of the Angel Prohibiting.}
\label{sermon:084}
\markright{Sermon 84. The Fact of St. John Is Declared, Which Would Have ...}

And I fell at his feet to worship him. And he said to me, “Do not do that. For I am your fellow servant, and one of your brothers, and of those who have the testimony of Jesus. Worship God. For the testimony of Jesus is the spirit of prophecy.”

\index{Augustine, Saint}\index{Daniel (Book of)}Here is the third point of this chapter, namely the actions of the Apostle John and the Angel of God. John wished to worship the Angel; but the Angel forbade him, instructing him to worship God. And before this act and undertaking of John seems chiefly to be considered. Angels are certainly noble creatures and of great power, by whom the Lord executes his greatest affairs. They often take on the form of men, and frequently appear to men, serve, protect, and do good to them, according to how God uses their ministry. For the Apostle, speaking of Angels (as I told you in Sermon 29), says they are not all ministering spirits, sent forth to serve those who will be heirs of salvation? And Scripture makes these things plain through various examples. Three appeared to Abraham in human likeness, who were Angels, instructing him; two delivered Lot himself out of the hands of the Sodomites and brought him out of the fire; whole armies of Angels surrounded Jacob, defending him against the force and violence of his brother Esau. The Lord sent his Angel before Moses and the children of Israel to lead them through the wilderness into the land of promise. Fiery chariots surrounded Elisha. An Angel lifted the siege of Jerusalem, slaying one hundred eighty-five thousand Assyrians. Daniel had Angels familiar with him. Likewise, the fathers and other Prophets. An Angel delivers Joseph out of all care; the same delivers the wise men from Herod’s treachery; immediately he commands that Christ be conveyed away to Egypt. Angels minister to Christ; in white garments they testified that the Lord was risen and ascended into heaven. The same brought the Apostles out of prison; one of them delivered Peter out of Herod’s prison. An Angel was sent to Cornelius, an Italian captain. Angels often spoke with Paul. They frequently bestow great benefits upon men. They declare themselves through God to be of great power. And while people observe these things, they would worship Angels—as is even the case now, where the Apostle John understood that Christ himself, by his Angel, had opened to him such great mysteries for the profit of the churches. While he marveled at his brightness and godly gifts, he was about to worship this Angel, the bringer of mysteries. Not that he intended or purposed to revolt from God, and to worship an Angel instead of God—for neither is it lawful to imagine such wickedness of so great an Apostle. He would therefore have worshipped and honored the Angel with \textit{Dulia}, as they term it (and as Augustine explains it), not with \textit{Latria}: that is to say, to worship and honor God as God, but the Angel somewhat less, as an excellent messenger of God. However, in this he offended, so that all men should understand that they sin, however many worship and honor Angels or excellent creatures with godly worship. As all the worshippers of Saints do at this day in Papistry. Neither do they have any other way to excuse their error than that same distinction, that God is worshipped and honored with worship \textit{latrical}, and Angels and Saints with worship \textit{dulical}, and the Virgin Mary with honor \textit{hiperdulical}, and I know not what other things, which I am both ashamed and loath to rehearse.

\index{John, Saint}And it appears that St. John here was entangled with the same error, whom otherwise we must needs confess to have sinned by Apostasy, and would have worshipped the Angel for God, or with God. Which are both two wicked, and unworthy such a man. But if he worshipped God, and nevertheless would have worshipped the Angel also, what thing else did he do than offend in the worship itself? And verily God hath permitted so worthy a man to err, as he did also Peter and Thomas; to the intent he might heal their infirmities: that is to wit, that by their errors we might learn to believe more rightly, and to honor God more purely. For this present place teaches openly, and other like examples of errors, that all the sayings and doings of Saints are not to be allowed without examination.

Now, consider the matter of a most excellent angel, that is to say, a godly refutation of error. First, he does not lightly say, “Do not as you have planned,” but severely condemning his action, he says with a certain vehemence, “See that you do it not.” We have a similar phrase in Switzerland, when (signifying that one must be very careful) we say, “Long und thu das nitte.” See that you do it not. Therefore, we have learned by the testimony of the angel that now neither angels nor saints are to be worshipped. For the Lord himself says of saints, “They shall be as the angels of God.” I do not see why they should not be equated with angels. And we have certainly learned that they may be worshipped neither with liturgical nor secular worship. And to worship is to honor with the mind, to fall at the feet, to bow down and kneel, as I have said elsewhere.

\index{Antichrist}After the angel shows reasons why he ought not to worship, for I am your fellow servant. He does not say “servant,” but “fellow servant,” meaning of the same office as you, under the same lord and master. For angels serve God after their manner, and so do men serve God after their manner; yet all are servants, and that of one master. And it is against reason that one servant should honor and worship another of his fellows, being of the same state and creation. It is therefore an unworthy matter that the faithful should worship the Apostles, Prophets, or Martyrs; much less does it become them to honor their dead bones. And lest any man should say, how the angel, indeed in respect of the most excellent Apostle John, confesses himself to be his fellow servant; but that there is another consideration to be had of other men, who do not come near the dignity of blessed John: and therefore we are much inferior, we may worship angels and apostles our superiors, he prevents and says, “and of your brethren.” And who are the brethren of the Apostle John? The angel himself answers, and says, “which have the testimony of Jesus.” The testimony of Jesus is the gospel, and the very faith fixed on the gospel, comprehending with a faithful mind Jesus. Wherefore all the faithful of Christ are John’s brethren; therefore the angel is their fellow servant also. And therefore none of the faithful ought to worship any angel or apostle; the Lord himself also in Matthew 12 calls all that obey his word or preaching brethren. And here is diligently to be noted, that by faith we are made the brethren of Christ, of angels and apostles. This the monks and friars should have beaten in and set forth, and not the fraternity of our Lady and fraternities of Saints, unless they had been the Apostles of that great and abominable Antichrist.

Moreover, the angel himself, explaining his own words again, shows what the testimony of Jesus Christ is. For the testimony of Jesus is the spirit of prophecy. And the spirit signifies revelation or understanding, and prophecy the prophetical and apostolic doctrine. Therefore, the meaning is: the testimony of Jesus Christ is nothing other than the revealing of the doctrine of prophets and apostles in the mind of the faithful through the Holy Spirit and faith. And therefore, the apostles are called witnesses in the gospel, and the gospel is a testimony. To testify is to preach. From this explanation, the following argument may be drawn: the reason for your worshipping me, John, is undoubtedly that excellent revelation and prophecy that I have revealed to you. But if I should therefore seem worthy of worship because there is in me an excellent spirit of prophecy, then by the same reasoning you should worship all your brethren, in whom is the same spirit of prophecy, namely, the testimony of Jesus, the true faith. But when you see this, and are compelled to grant it, you will find it very absurd. I perceive it to be absurd if you should worship an angel.

The last and strongest reason why he would not be worshipped is this: worship God. This command is derived from God’s eternal and unchanging authority and law, revealed in Deuteronomy 6 and repeated by our Savior Christ in Matthew 4. If we would obey God’s law, all cults, worship, and invocations of saints would long ago have been banished and exiled from the church.

Furthermore, there are other passages that praise the ministries and virtues of angels, yet teach that we should honor and call upon God himself. Read the excellent Psalms 34 and 91. And if anyone desires the agreement of the church fathers, let him read Augustine, saying that angels must neither be worshipped nor called upon, nor should sacrifices be made to them, nor churches erected. The key passages are found in the writings on \textit{True Religion}, chapter 55, \textit{Against Maximinus, an Arian Bishop}, first book, page 77; \textit{The City of God}, book 8, last chapter; chapters 10, 16, 19, and 20. To God be the glory.
